\section{Background}

\subsection{Agent Lifecycle Management in Multi-Agent Systems}

Multi-agent systems (MAS) have evolved from simple reactive architectures~\cite{brooks1991intelligence} to sophisticated frameworks capable of autonomous reasoning, collaboration, and problem-solving at scale. A foundational challenge in MAS research is the \emph{agent lifecycle management problem}: how to create, initialize, coordinate, monitor, and terminate agents in a manner that preserves system integrity, enables compositional reasoning, and supports fault isolation~\cite{horling2004survey,wooldridge2009multiagent}. Traditional agent platforms---including Jade~\cite{bellifemine2007jade}, SPADE~\cite{palanque2012spade}, and Microsoft AutoGen~\cite{wu2023autogen}---treat agent creation as a stateless, ephemeral event. An agent is instantiated, performs its designated task, and is discarded. There is no inherent mechanism to record the identity of the spawning entity, the state at spawn time, or the chain of delegation that led to a given computation. This omission is not merely a bookkeeping deficiency; it fundamentally undermines accountability, debuggability, and regulatory compliance in production deployments~\cite{singh2019blockchain}.

Recent advances in large language model (LLM)-based agents have intensified this problem. Systems such as AutoGPT~\cite{richards2023autogen}, LangChain Agents~\cite{chase2022langchain}, and MetaGPT~\cite{hong2023metagpt} construct pipelines of specialized agents that collaborate to solve complex tasks. These systems routinely spawn sub-agents, delegate objectives, and synthesize results, but they typically do so without preserving a durable provenance graph. When a system's output is incorrect or its reasoning opaque, practitioners lack the forensic substrate to reconstruct \emph{why} a particular agent arrived at a given conclusion. The problem is not one of implementation detail; it is architectural. Unless the framework itself encodes the parent-child relationship as a first-class, immutable construct, it will remain a matter of convention rather than constraint.

The AgentOS framework introduced by Beebe and Blanco~\cite{beebe2025agentos} recognizes this as a core design failure of existing agent platforms. AgentOS proposes that every agent's origin, delegation chain, and execution context be recorded as an immutable append-only log attached to the agent's identity. This positions lifecycle management not as a deployment concern but as a structural property of the agent substrate itself.

\subsection{Accountability and Provenance in Distributed Systems}

The problem of accountability in distributed computation has been studied extensively in the context of cloud computing, microservices, and blockchain. A general consensus has emerged: accountability requires that every action in a system be attributable to an identified principal, and that the causal chain linking inputs to outputs be inspectable~\cite{weitzner2008information,bates2015accountable}. In blockchain systems, this is achieved through transaction logs and Merkle proofs~\cite{nakamoto2008bitcoin}. In cloud auditing frameworks such as AWS CloudTrail~\cite{awscloudtrail} and Google Cloud Audit Logs~\cite{googleauditlogs}, it is achieved through structured event logging. In both cases, the underlying principle is the same: a tamper-evident record of causality that can be independently verified.

Applying this principle to AI agent systems introduces unique challenges. Unlike a REST API call or a database transaction, an agent's computation is nondeterministic, context-dependent, and often semi-structured. The notion of a ``principal'' is complicated by the fact that an agent may act on behalf of a user, on behalf of another agent, or on behalf of a organizational role. Nevertheless, the core requirement is unchanged: there must exist an authoritative, immutable record of the delegation chain that caused a particular computation to occur.

Prior work on provenance in scientific workflows~\cite{buneman2001 provenance,frew2008reproducible} established that capturing the \emph{dataflow graph} of a computation is essential for reproducibility and attribution. Extensions to cloud workflows~\cite{simmhan2005survey} formalized provenance as a first-class citizen. More recently, the AI ethics literature has argued that autonomous systems must maintain ``audit trails'' analogous to those in financial systems, particularly when agents make decisions that affect individuals~\cite{buvac2021artificial,brundage2020toward}. The BX3 Framework extends this lineage thinking into the agent substrate layer, treating the delegation chain not as an optional metadata field but as the structural backbone of agent identity.

\subsection{Fixed vs. Mutable Delegation Chains}

A critical distinction in delegation architecture is whether the chain of custody---the sequence of principals through which authority is transmitted---is \emph{fixed} at creation time or \emph{mutable} after the fact. In traditional access control models, delegation is often mutable: a principal can re-delegate its authority, revoke previously granted permissions, or modify the delegation graph post hoc~\cite{ellison1999 SPKI,later2002pem}. While flexible, mutable delegation introduces significant risks. It becomes possible for a delegation chain to be altered retrospectively, obscuring the actual sequence of decisions that led to an outcome. This is particularly problematic in adversarial settings, where actors may have incentives to retroactively construct favorable narratives of authorization.

Immutable delegation chains offer a different tradeoff. Once a principal $A$ delegates authority to principal $B$, the delegation record is sealed. $B$ cannot alter its own lineage; it can only further delegate to $C$, appending to the chain rather than modifying it. This model is well-established in certificate transparency logs~\cite{laurie2014certificate} and in the delegated proof-of-stake mechanisms of modern blockchain systems~\cite{king2012ppcoin}. It provides strong non-repudiation: any observer can verify the exact sequence of delegations that preceded a given action.

The BX3 Framework adopts the immutable delegation model as a foundational design principle. In BX3, every agent is assigned a \emph{parent chain} at creation: a cryptographically linked sequence of the agents that preceded it in the delegation hierarchy. This chain is append-only. An agent may spawn a child agent, but it cannot modify its own parent chain. The resulting structure is a growing tree (or forest, at the system level) of agent lineages, each node annotated with its spawn context, model configuration, and execution state at the time of delegation. This approach provides strong provenance guarantees without sacrificing the flexibility of hierarchical decomposition.

Mutable delegation, by contrast, remains the norm in existing LLM agent frameworks. In AutoGen, an agent can hand off control to another agent, but the handoff is a transient message-passing event with no durable record of the delegation relationship. In LangChain, tool use and agent chains are defined at runtime but are not attached to the agent's identity in a way that survives process termination. These systems optimize for flexibility and dynamic reconfiguration, but they sacrifice the forensic guarantees that production deployments demand.

\subsection{Hierarchical Task Decomposition in AI}

Hierarchical task decomposition is a foundational technique in classical AI planning, robotics, and problem-solving systems. The basic insight, articulated in seminal work by Sacerdoti~\cite{sacerdoti1977hierarchy} and later refined in the Non-Linear Planning framework~\cite{chapman1987planning}, is that complex goals can be productively decomposed into subgoals, which can themselves be further decomposed until they reach primitive actions that can be executed directly. This approach mitigates the combinatorial explosion that plagues flat search-based planning by imposing a structural bias toward shallow, branching trees of increasing concreteness.

In the context of modern LLM-based agents, hierarchical decomposition has re-emerged as a dominant design pattern. The \emph{Hierarchical Planning and Acting} (HPA) architecture~\cite{planchet2024hierarchical} formalizes the intuition that an LLM acting as a top-level planner can generate a high-level task graph, while specialized sub-agents handle individual task nodes. Similarly, the \emph{Tree of Thoughts} framework~\cite{yao2023tree} and its successors~\cite{braese2024process} explore spaces of reasoning by maintaining a branching structure of partial solutions, guided by an orchestrator that evaluates and prunes branches. MetaGPT~\cite{hong2023metagpt} implements a structured multi-agent workflow in which roles are defined hierarchically and communicate via a shared message protocol, mimicking an organizational chart.

The key limitation of these approaches, from a provenance perspective, is that the decomposition tree exists only at the level of \emph{control flow}, not at the level of \emph{agent identity}. When a parent agent decomposes a task and delegates sub-tasks to child agents, the resulting child processes are not represented as first-class objects with their own immutable lineage. Their existence is ephemeral; they are created, used, and discarded within the scope of a single planning cycle. The decomposition structure is implicit in the control flow, not encoded as a durable data structure.

The BX3 Framework addresses this by treating each decomposition step as an agent spawn event. When an agent decomposes a task, it formally spawns child agents, each of which receives an immutable parent chain recording the decomposition lineage. The task itself becomes a first-class artifact of the agent ecosystem, with its own identity, status, and associated provenance graph. This transforms hierarchical decomposition from a transient planning heuristic into a persistent, auditable structure that can be inspected, explained, and verified after the fact.

\subsection{The BX3 Framework as Substrate for Immutable Parent Chains}

The BX3 (Broadcast $\times$ Execute $\times$ eXchange) Framework, introduced by Beebe and Blanco~\cite{beebe2025bx3}, is a domain-agnostic substrate for building, composing, and reasoning about autonomous agent systems. BX3 departs from conventional agent frameworks in several fundamental respects, most notably in its treatment of agent identity and delegation. Where existing frameworks treat an agent as an isolated process with transient communication channels, BX3 treats every agent as a node in a persistent, immutable lineage graph. This graph is not a metadata annotation; it is the primary data structure around which the entire system is organized.

The core abstraction in BX3 is the \emph{Agent Chain}, a cryptographically signed, append-only sequence that records the spawning of each agent. Formally, an agent $A_n$ is defined as a tuple:
\begin{equation}
A_n = \langle \text{id}_n,\ \text{parent\_chain}_n,\ \text{context}_n,\ \text{model}_n,\ \text{state}_n \rangle
\end{equation}
where $\text{id}_n$ is a unique identifier, $\text{parent\_chain}_n = [A_0, A_1, \dots, A_{n-1}]$ is the immutable sequence of ancestors, $\text{context}_n$ captures the spawn-time execution context, $\text{model}_n$ specifies the model configuration, and $\text{state}_n$ is the mutable runtime state. The parent chain $\text{parent\_chain}_n$ is set at agent creation and cannot be modified. When $A_n$ spawns a child agent $A_{n+1}$, the resulting $\text{parent\_chain}_{n+1}$ is defined as $\text{parent\_chain}_n \oplus [A_n]$, where $\oplus$ denotes concatenation. This ensures that the full ancestral lineage is always recoverable from any agent in the system.

BX3 further defines a \emph{Recursive Spawning Protocol} (RSP) that governs the conditions and mechanics of agent delegation. RSP specifies that any agent may spawn child agents to handle sub-tasks, but that every spawn event must record the delegation context, the resource budget allocated to the child, and the termination criteria. These records are embedded in the child agent's $\text{context}$ field, providing a complete audit trail of \emph{why} a given sub-task was delegated, \emph{to whom}, and \emph{with what constraints}. This stands in contrast to conventional agent frameworks, where delegation is implicit in message passing and leaves no structural trace.

The implications of this design are significant. Because every agent carries its complete parent chain, any output produced by the system can be traced back through the exact sequence of decompositions, delegations, and computations that produced it. This enables what BX3 terms \emph{chain-of-thought provenance}: not merely the reasoning steps of a single model, but the full hierarchical history of which agent decomposed which task, which sub-agent executed which component, and how the results were synthesized. This property is essential for compliance with emerging AI governance frameworks that require explainability and auditability of autonomous systems~\cite{eu2024aiact, NIST-AI-2023}.

The substrate is designed to be model-agnostic. BX3 does not prescribe a specific LLM backend; instead, it provides a common interface through which any compatible inference engine can be plugged in. This ensures that the immutable parent chain mechanism is orthogonal to the choice of underlying model, making the framework applicable across the rapidly evolving landscape of foundation model providers. Preliminary implementations have demonstrated RSP compatibility with Ollama-based local inference~\cite{ollama2024} and with cloud-based API providers including Mistral AI~\cite{mistral2024}, providing evidence that the framework is viable across deployment contexts.
